\subsection{Lemma contracts and remaining comparisons}

The first contract is a finite-abelian corollary with an explicit phase and
central-quotient comparison. The underlying tensor theorem is already
available. The statements below distinguish these inherited results from
the remaining affine naturality, relation semantics and analytic work.
The first nine propositions are proved; the remaining six retain sketch status.

\begin{proposition}[Weyl realization]
\statusProved{}
For every odd-characteristic finite field $k$ of cardinality $q$, nontrivial additive
character $\psi$, and symplectic $k$-space $V$ of dimension $2n$, at
$\beta_V=\omega_V/2$ the Weyl algebra is a unital $*$-algebra with faithful
normalized trace $\tau_V$ and is $*$-isomorphic to $\operatorname{End}(\mathbb
C^{q^n})$. Its standard formula is a unitary irreducible realization;
irreducible unitary representations of the prescribed Heisenberg group with central character $\psi$ are unique up to unitary equivalence. Unitary
intertwiners between irreducible unitary models are unique up to $U(1)$. Rank zero
gives $\mathbb C$; rank one specializes to the existing Weyl algebra at
$\beta=\omega/2$.
\end{proposition}
\begin{scope}
Odd characteristic only; no preferred abstract-space basis or genuine symmetry lift is asserted.
\end{scope}
\provenance{SP-WEYL}{Structured proof: reuse.md, sections 1--2; inherited results as cited}{Exact finite probes at their declared scope}{Single blind review and stage-1 adjudication}

\begin{proof}
If $V$ is nonzero, choose a nonzero vector $e\in V$.  Nondegeneracy
supplies $f$ with $\omega_V(e,f)=1$, and the plane spanned by $e,f$ is
nondegenerate.  Its orthogonal complement is again symplectic, so induction
gives a symplectic coordinate isomorphism
\[
 V\simeq k^n\oplus k^n,
 \qquad
 \omega_V((a,b),(a',b'))=a\mathbin{\cdot}b'-a'\mathbin{\cdot}b.
\]
For $V=0$, use the unique coordinate map between zero spaces; this case is
also recovered below from the empty configuration group.
The image of $\psi$ is the group of $p$th roots selected by $\psi$.  For
$A=(k^n,+)$, the characters
\[
 \chi_b(x)=\psi(-b\mathbin{\cdot}x)
\]
are all the characters of $A$: if $b\ne0$, the functional
$x\mapsto b\mathbin{\cdot}x$ is onto, so $\chi_b$ is nontrivial, and the
previously proved finite-character count turns this injection into a
bijection.

In these coordinates the finite-abelian Weyl multiplier is
$\psi(a\mathbin{\cdot}b')$.  Rephase its operators by
\[
 c(a,b)=\psi(a\mathbin{\cdot}b/2),
 \qquad
 W^{\mathrm s}(a,b)=c(a,b)W_{\mathrm{ref}}(a,\chi_b).
\]
Expanding the dot product gives
\begin{align*}
 &a\mathbin{\cdot}b'
 +\frac12a\mathbin{\cdot}b
 +\frac12a'\mathbin{\cdot}b'
 -\frac12(a+a')\mathbin{\cdot}(b+b') \\
 &\hspace{35mm}
 =\frac12\bigl(a\mathbin{\cdot}b'-a'\mathbin{\cdot}b\bigr).
\end{align*}
Thus the rephased operators have exactly the half-form multiplier and satisfy
$(W^{\mathrm s}(v))^*=W^{\mathrm s}(-v)$.  On the computational basis,
\[
 W^{\mathrm s}(a,b)\delta_y
 =\psi\!\left(-b\mathbin{\cdot}y-\frac12a\mathbin{\cdot}b\right)
   \delta_{y+a},
\]
which is the stated wavefunction formula after writing the output coordinate
as $x=y+a$.

Rephasing by nonzero scalars does not change the operator span.  The finite
Stone--von Neumann theorem already proved for $A$ therefore gives all of
$\operatorname{End}(\mathbb C^A)$, with
\[
 \operatorname{Tr}(W^{\mathrm s}(v))=q^n\,\mathbf 1_{v=0},
 \qquad
 \tau_V=q^{-n}\operatorname{Tr}.
\]
This trace is normalized and faithful because
$\operatorname{Tr}(B^*B)>0$ for every nonzero matrix $B$.

It remains to compare central coordinates.  The homomorphism
\[
 (t,a,b)\longmapsto
 \bigl(\psi(t+a\mathbin{\cdot}b/2),a,\chi_b\bigr)
\]
from the prescribed $k$-central Heisenberg group onto the finite-abelian
phase group has kernel
$\{(t,0,0):\psi(t)=1\}$.  A unitary representation with the stipulated
central character kills this kernel, and conversely a unitary representation
of the quotient pulls back with that central character.  The previously
proved finite Stone--von Neumann result now supplies dimension $q^n$,
irreducibility, unitary equivalence, and uniqueness of unitary intertwiners
up to $U(1)$.  For $n=0$ the construction is $\mathbb C$; for $n=1$ it is
the already defined rank-one half-form Weyl algebra.
\end{proof}

\begin{proposition}[Affine symmetries and projective implementation]
\statusProved{}
For a fixed odd-characteristic finite field $k$ and nontrivial additive character
$\psi$, put $\beta_V=\omega_V/2$ on each symplectic space. An affine symplectic arrow
$(t,g):V\to W$ induces the $*$-isomorphism
$\alpha_{t,g}(W_{\beta_V}(v))=\psi(\omega_W(t,gv))W_{\beta_W}(gv)$. These assignments
respect identities and affine composition exactly. On irreducible unitary models they
have unique projective unitary implementers. Composition of those implementers agrees
in the $U(1)$ quotient.
\end{proposition}
\begin{scope}
This is a covariant algebra functor and a projective implementation statement. Choosing a genuine linear lift is a separate comparison.
\end{scope}
\provenance{SP-EGOROV}{Structured proof: egorov.md; inherited results as cited}{Exact finite probes at their declared scope}{Single blind review and stage-1 adjudication}

\begin{proof}
It is enough first to check the prescribed map on the Weyl basis.  For
$v,z\in V$, symplecticity of $g$ gives
\begin{align*}
 \alpha_{t,g}(W_{\beta_V}(v))\alpha_{t,g}(W_{\beta_V}(z))
 &=\psi\!\left(
      \omega_W(t,gv)+\omega_W(t,gz)+\frac12\omega_W(gv,gz)
    \right)W_{\beta_W}(g(v+z)) \\
 &=\psi\!\left(
      \omega_W(t,g(v+z))+\frac12\omega_V(v,z)
    \right)W_{\beta_W}(g(v+z)) \\
 &=\alpha_{t,g}
   \bigl(W_{\beta_V}(v)W_{\beta_V}(z)\bigr).
\end{align*}
The unit is fixed. Since $\overline{\psi(x)}=\psi(-x)$, the same basis gives
\[
 \alpha_{t,g}(W_{\beta_V}(v)^*)=
 \alpha_{t,g}(W_{\beta_V}(v))^*.
\]
The inverse affine arrow is
$(-g^{-1}t,g^{-1})$, so this homomorphism is a $*$-isomorphism.

The semidirect-product law is exact, rather than projective.  If
$(t,g):V\to W$ and $(s,h):W\to Z$, then
\begin{align*}
 (\alpha_{s,h}\alpha_{t,g})(W_{\beta_V}(v))
 &=\psi\!\left(\omega_W(t,gv)+\omega_Z(s,hgv)\right)
   W_{\beta_Z}(hgv) \\
 &=\psi\!\left(\omega_Z(s+ht,hgv)\right)W_{\beta_Z}(hgv) \\
 &=\alpha_{s+ht,hg}(W_{\beta_V}(v)).
\end{align*}
The identity case follows by taking $(t,g)=(0,1_V)$.  Pure translation is
inner:
\[
 W_{\beta_W}(t)W_{\beta_W}(v)W_{\beta_W}(t)^*
 =\psi(\omega_W(t,v))W_{\beta_W}(v),
\]
so the affine formula is the Weyl translation followed by the linear
Egorov map.

Let $(H_V,\pi_V)$ and $(H_W,\pi_W)$ be objects of the arbitrary-rank model
groupoids just defined.  Pulling $\pi_W$ back along $\alpha_{t,g}$ gives
another object of the source model groupoid.  Finite Stone--von Neumann
uniqueness supplies a unitary
$U_{t,g}:H_V\to H_W$ with
\[
 U_{t,g}\pi_V(a)U_{t,g}^*=\pi_W(\alpha_{t,g}(a)),
 \qquad a\in A_{\psi,\beta_V}(V),
\]
and any two such unitaries differ by a phase.  For two composable affine
arrows, $U_{s,h}U_{t,g}$ and $U_{s+ht,hg}$ implement the same exact algebra
map calculated above.  Their classes therefore agree after quotienting by
$U(1)$.  This proves projective composition without selecting a genuine
phase lift.
\end{proof}

\begin{proposition}[Symplectic sum and quantum tensor compatibility]
\statusProved{}
For a fixed odd-characteristic finite field $k$ and nontrivial additive character
$\psi$, put $\beta_V=\omega_V/2$, $\beta_W=\omega_W/2$. The assignment
$W_{\beta_V\oplus\beta_W}(v,w)\mapsto W_{\beta_V}(v)\otimes W_{\beta_W}(w)$ gives a
natural trace-preserving $*$-isomorphism $A_{\psi,\beta_V\oplus\beta_W}(V\oplus W)\to
A_{\psi,\beta_V}(V)\otimes A_{\psi,\beta_W}(W)$. Together with the rank-zero unit
these maps obey the associativity, unit and symmetry diagrams. The corresponding
unitary model identifications obey these diagrams projectively.
\end{proposition}
\begin{scope}
The source operation is symplectic direct sum. This does not construct a coherent additive completion of the classical category.
\end{scope}
\provenance{SP-TENSOR}{Structured proof: tensor.md; inherited results as cited}{Exact finite probes at their declared scope}{Single blind review and stage-1 adjudication}

\begin{proof}
Temporarily denote the displayed generator map by $\Theta_{V,W}$.  The
direct-sum form gives
\begin{align*}
 &\Theta_{V,W}(W_{\beta_V\oplus\beta_W}(v,w))
  \Theta_{V,W}(W_{\beta_V\oplus\beta_W}(v',w')) \\
 &\quad=
 \psi\!\left(\frac12\omega_V(v,v')+\frac12\omega_W(w,w')\right)
 W_{\beta_V}(v+v')\otimes W_{\beta_W}(w+w'),
\end{align*}
which is the image of the product in the source algebra.  The map also
preserves the unit and involution.  It takes the Weyl basis indexed by
$V\oplus W$ bijectively to the tensor-product Weyl basis, and hence is a
$*$-isomorphism.  On these basis elements,
\[
 (\tau_V\otimes\tau_W)
 \bigl(W_{\beta_V}(v)\otimes W_{\beta_W}(w)\bigr)
 =\mathbf 1_{v=0}\mathbf 1_{w=0}
 =\tau_{V\oplus W}(W_{\beta_V\oplus\beta_W}(v,w)),
\]
so it preserves the normalized trace.

In standard coordinates, reorder
$((a_1,b_1),(a_2,b_2))$ as
$((a_1,a_2),(b_1,b_2))$.  The previously proved configuration-product
unitary sends $\delta_{(x,y)}$ to $\delta_x\otimes\delta_y$.  The half-form
phase is compatible with this unitary because
\[
 \psi\!\left(\frac12(a_1\mathbin{\cdot}b_1
                         +a_2\mathbin{\cdot}b_2)\right)
 =\psi(a_1\mathbin{\cdot}b_1/2)\,
  \psi(a_2\mathbin{\cdot}b_2/2).
\]
Thus the standard Weyl operator on the direct sum is carried exactly to the
tensor of the two standard Weyl operators.

Naturality for translations and linear symplectic maps is also exact.  For
affine arrows $(t,g):V\to V'$ and $(s,h):W\to W'$, both routes around the
naturality square send the generator labelled by $(v,w)$ to
\[
 \psi\!\left(\omega_{V'}(t,gv)+\omega_{W'}(s,hw)\right)
 W_{\beta_{V'}}(gv)\otimes W_{\beta_{W'}}(hw).
\]
Rebracketing three summands, adjoining the zero space, or exchanging two
summands likewise sends every Weyl generator to the same ordered tensor by
either route.  The associativity, unit and symmetry diagrams therefore
commute at the algebra level.

For arbitrary objects $(H_V,\pi_V)$, $(H_W,\pi_W)$, and
$(H_{V\oplus W},\pi_{V\oplus W})$ of the model groupoids just defined,
finite Stone--von Neumann uniqueness supplies an implementing unitary
\[
 J_{V,W}:H_{V\oplus W}\longrightarrow H_V\otimes H_W
\]
for the displayed algebra isomorphism, unique up to $U(1)$.  In an affine
naturality diagram, the two model-level routes implement the same exact
algebra map just calculated, so they differ by one phase.  The same argument
applied to the exact algebra associativity, unit and symmetry diagrams shows
that every corresponding model diagram commutes in the $U(1)$ quotient.
This uses only projective classes and makes no coherent choice of unitary
representatives.
\end{proof}

\begin{proposition}[Composition of affine Lagrangian relations]
\statusProved{}
For every finite field $k$, the affine Lagrangian relation prescription is closed
under relational composition and product and forms a dagger symmetric monoidal
category. Its unit is the zero symplectic space. Graphs of affine symplectic
isomorphisms define a faithful symmetric monoidal functor from the affine symmetry
groupoid.
\end{proposition}
\begin{scope}
This concerns finite linear/affine geometry, including empty relations. It does not assert a normalized quantization of those relations.
\end{scope}
\provenance{SP-LREL}{Reviewed structured relation proofs}{Exact F2/F3 relation and compact controls}{Single blind review, sole repair and relation-cluster adjudication}

\begin{proof}
For linear arrows $R\subseteq\overline V\oplus W$ and
$S\subseteq\overline W\oplus Z$, put
\[
 X=\overline V\oplus W\oplus\overline W\oplus Z,
 \quad L=R\oplus S,
 \quad C=\overline V\oplus\Delta_W\oplus Z.
\]
Then $L$ is Lagrangian, $C$ is coisotropic, and
\[
 C^\perp=\{(0,w,w,0):w\in W\}.
\]
The reduced space $C/C^\perp$ is symplectically identified with
$\overline V\oplus Z$ by
$[(v,w,w,z)]\mapsto(v,z)$.  For completeness, if $K=C^\perp$,
$\dim X=2N$ and $\dim K=r$, then
\[
 (L+K)^\perp=L\cap C,
 \qquad
 \dim(L\cap C)=N-r+\dim(L\cap K).
\]
It follows that the image of $L\cap C$ in $C/K$ is isotropic of dimension
$N-r$, half the reduced dimension, hence is Lagrangian.  Under the displayed
identification this image is exactly
\[
 \{(v,z):\exists w,\ (v,w)\in R,\ (w,z)\in S\}=S\circ R.
\]
For nonempty affine translates, intersection with $C$ is either empty or,
after choosing one point $(v_0,w_0,w_0,z_0)$, is that point plus the
intersection of the direction space with $C$.  Its outer projection is
therefore $(v_0,z_0)$ plus the linear composite.  This proves affine and
empty closure without a transversality assumption and without division, so
it includes characteristic two.

Associativity and the identity laws follow by expanding the existential
middle variables.  Converse exchanges the two endpoint factors; this map
negates the ambient form and hence preserves Lagrangian subspaces, including
in characteristic two.  Direct sum preserves Lagrangian direction spaces,
and the prescribed reorder is symplectic because
\[
 (-\omega_V+\omega_W)+(-\omega_{V'}+\omega_{W'})
 =-(\omega_V+\omega_{V'})+(\omega_W+\omega_{W'}).
\]
The usual rebracketing, zero-space and exchange graphs give the coherence
maps, whose diagrams reduce to equality of tuple maps.  Finally, for an
affine symplectic isomorphism $f(v)=gv+t$, its graph is the translate of
$\{(v,gv)\}$; the latter is isotropic of half dimension.  Graphs preserve
identity, composition and direct sum and determine their functions, giving
the faithful symmetric monoidal functor.
\end{proof}

\begin{proposition}[Dagger compact affine Lagrangian relations]
\statusProved{}
For every finite field $k$, the prescribed compact data make
$\mathsf L_k^{\mathrm{aff}}$ a dagger compact category with dual $\overline V$.
The diagonal cup and cap are affine Lagrangian relations and satisfy
\[
\begin{aligned}
 \lambda_V\circ(\epsilon_V\oplus1_V)\circ a^{-1}_{V,\overline V,V}
 \circ(1_V\oplus\eta_V)\circ\rho_V^{-1}&=1_V,\\
 \rho_{\overline V}\circ(1_{\overline V}\oplus\epsilon_V)
 \circ a_{\overline V,V,\overline V}\circ(\eta_V\oplus1_{\overline V})
 \circ\lambda_{\overline V}^{-1}&=1_{\overline V}.
\end{aligned}
\]
They obey $\eta_V^\dagger=\epsilon_V\circ\sigma_{\overline V,V}$.
Name and unname are inverse bijections
$\mathsf L_k^{\mathrm{aff}}(V,W)\cong
\mathsf L_k^{\mathrm{aff}}(0,\overline V\oplus W)$ that retain the same
ordered affine relation and preserve the empty relation.
The endomorphisms of $0$ are exactly $\varnothing$ and $\Delta_0$.
Composition and tensor transported by $\lambda_0$ have the same
multiplication, with $\varnothing$ absorbing and $\Delta_0$ the identity.
For every $V$, $\epsilon_V\circ\sigma_{\overline V,V}\circ\eta_V=\Delta_0$.
\end{proposition}
\begin{scope}
Classical affine Lagrangian relations over every finite field, including characteristic two and empty relations. No quantum normalization, amplitude, probability, stabilizer equivalence or Choi theorem.
\end{scope}
\provenance{SP-COMPACT}{Reviewed structured relation proofs}{Exact F2/F3 relation and compact controls}{Single blind review, sole repair and relation-cluster adjudication}

\begin{proof}
The diagonal in $\overline V\oplus V$ is isotropic because its form is
$-\omega_V(v,w)+\omega_V(v,w)=0$ and is Lagrangian because its dimension is
$\dim V$.  This types both $\eta_V$ and $\epsilon_V$ with their displayed
factor orders.  Expanding the first snake, an input $x$ is related through
the cup to $(x,(u,u))$, the inverse associator sends this to $((x,u),u)$,
and the cap forces $x=u$; the final output is $x$.  Expanding the second
snake similarly sends $x$ through $((u,u),x)$ and the cap forces $u=x$.
Thus the two fully unital composites in the proposition are the relevant
identity relations.  Converse turns the cup into the diagonal effect on
$\overline V\oplus V$, which is
$\epsilon_V\circ\sigma_{\overline V,V}$.

For any $R:V\to W$, its name contains $(0,(v,w))$ exactly when
$(v,w)\in R$.  Conversely, the unname of $Q$ relates $v$ to $w$ exactly
when $(0,(v,w))\in Q$: the cap forces the copied input coordinate to equal
the first coordinate of the state.  Hence the two maps are inverse and
retain affine, nonfunctional and empty relations.

A relation $0\to0$ is a subset of a singleton, so the only two allowed
scalars are $\varnothing$ and $\Delta_0$.  Direct expansion shows that
composition and
$\lambda_0\circ(s\oplus t)\circ\lambda_0^{-1}$ have the two-element
Boolean multiplication table.  In the closed cup--cap composite the zero
vector is always a middle witness, so existential relation composition
gives $\Delta_0$ for every $V$; no cardinality or amplitude is retained.
\end{proof}

\begin{proposition}[The scalar-quotient stabilizer comparison]
\statusProved{}
For an odd prime $p$, restrict the affine Lagrangian relation category to the standard spaces $V_n=\mathbb F_p^n\oplus\mathbb F_p^n$, $n\geq0$. For every nonempty arrow $R:V_m\to V_n$, $\mathcal E_p(R)$ is a one-dimensional subspace of $\operatorname{Hom}_{\mathbb C}(H_{\mathbb F_p,m},H_{\mathbb F_p,n})$ contained in $\mathsf{Stab}^{\mathrm{amp}}_p(H_{\mathbb F_p,m},H_{\mathbb F_p,n})$, while $\mathcal E_p(\varnothing)=\{0\}$. The assignments $V_n\mapsto H_{\mathbb F_p,n}$, $R\mapsto[T]$ for any nonzero $T\in\mathcal E_p(R)$, and $\varnothing\mapsto[0]$ define a dagger symmetric monoidal equivalence to $\mathsf{Stab}^{\mathrm{proj}}_p$.
\end{proposition}
\begin{scope}
Standard-object props over odd prime fields, modulo all invertible complex scalars, retaining a separate zero. No norm, success probability, CP map or arithmetic-source exhaustion.
\end{scope}
\provenance{SP-STAB-REL}{Reviewed intertwiner-line, functor and equivalence proofs}{Exact stabilizer comparison controls}{Single blind review, sole repair and relation-cluster adjudication}

\begin{proof}
On the Hilbert--Schmidt space
$\operatorname{Hom}_{\mathbb C}(H_{\mathbb F_p,m},H_{\mathbb F_p,n})$
put
\[
 \rho(v,w)T=W^{\mathrm s}_{\mathbb F_p,n}(w)T
 W^{\mathrm s}_{\mathbb F_p,m}(v)^*.
\]
The half-form Weyl law gives
\[
 \rho(x)\rho(y)=\psi_{\mathbb F_p}
   (\Omega_{m,n}(x,y)/2)\rho(x+y),
\]
and the trace of $\rho(x)$ is zero unless $x=0$, when it is
$p^{m+n}$.  If $R=r+L$ is nonempty, isotropy makes $\rho|_L$ an ordinary
unitary representation and
\[
 P_R=\frac1{|L|}\sum_{\ell\in L}
 \psi_{\mathbb F_p}(\Omega_{m,n}(r,\ell))\rho(\ell)
\]
is the orthogonal projector onto the displayed intertwiner space.  Changing
$r$ by an element of $L$ leaves every coefficient unchanged.  Since
$|L|=p^{m+n}$, Weyl trace orthogonality gives
$\operatorname{Tr}(P_R)=1$.  Thus every nonempty relation gives one nonzero
operator line, while the empty relation gives only zero.

Composition also detects empty affine intersections.  For
$R:V_m\to V_n$, let
\[
 A_R=\{a:(0,a)\in L_R\};
\]
the range of a nonzero $T\in\mathcal E_p(R)$ is exactly the simultaneous
Weyl eigenspace for $A_R$ with character
$a\mapsto\psi_{\mathbb F_p}(-\omega_n(r_n,a))$.  The analogous input
eigenspace describes the initial support.  For composable nonempty $R,S$,
the two middle support projectors have nonzero product exactly when their
characters agree on the intersection of their isotropic label groups.  By
symplectic annihilators, this is exactly the condition that the affine middle
projections meet, hence exactly that $S\circ R$ is nonempty.  In that case,
choosing a common middle origin cancels the two middle form terms and proves
\[
 \mathcal E_p(S)\mathcal E_p(R)=\mathcal E_p(S\circ R).
\]
When the middle projections do not meet, the projector product and every
operator composite are zero.

For bare relational converse the endpoint form changes sign.  Taking the
Hilbert adjoint of the defining equation gives exactly
\[
 \mathcal E_p(R^\dagger)=\mathcal E_p(R)^*.
\]
The direct-sum form adds and the admitted grouped Weyl comparison gives
$\mathcal E_p(R\oplus S)=\mathcal E_p(R)\otimes\mathcal E_p(S)$, including
the zero and rank-zero cases.  The identity relation has the scalar line of
the identity operator by the full-matrix Weyl commutant.

For target membership, use the explicit symplectic identification
$c_m:\overline V_m\to V_m$, $c_m(a,b)=(a,-b)$, and
\[
 \operatorname{vec}(T)=\sum_x\delta_x\otimes T\delta_x.
\]
Complex conjugation gives
$\overline{W^{\mathrm s}(a,b)}=W^{\mathrm s}(a,-b)$, so vectorization sends
$\mathcal E_p(R)$ to the stabilizer-state line of
$(c_m\oplus1)(\ulcorner R\urcorner)$.  Every affine Lagrangian state is a
Clifford image of $\delta_0^{\otimes(m+n)}$: extend a basis of its direction
to a symplectic basis and use the admitted affine Egorov implementer.
Unvectorization contracts with the stabilizer Bell effect, so every operator
in the line is an actual stabilizer amplitude.
For an empty arrow of type $V_m\to V_n$, let
$e_j:\mathbb C\to H_{\mathbb F_p,j}$ be the tensor of $j$ computational
zero preparations, with $e_0=1_{\mathbb C}$.  Its image is the typed actual
zero map
\[
 e_n\circ(0:\mathbb C\to\mathbb C)\circ e_m^*:
 H_{\mathbb F_p,m}\longrightarrow H_{\mathbb F_p,n},
\]
using only the preparation, scalar, tensor, adjoint and composition
generators.

Finally take an arbitrary $U\in C_2(A)$.  Conjugation has the unique form
\[
 U W(v)U^*=\gamma(v)W(gv).
\]
The Weyl product and commutator force $g$ to be linear symplectic and
$\gamma$ to be an additive character.  Hence uniquely
$\gamma(v)=\psi_{\mathbb F_p}(\omega(t,gv))$, so $U$ spans the line of the
affine graph $(t,g)$.  This treats every second-level unitary.  Together with
the computational preparation, its adjoint, and the scalar generators, it
proves fullness.  The projective Weyl stabilizer of a nonzero line recovers
its Lagrangian direction, and its eigencharacter recovers the affine coset,
which proves faithfulness.  The object assignment is surjective, completing
the stated dagger symmetric monoidal equivalence, using the relation and
compact laws proved above.
\end{proof}

\pagebreak[3]
\begin{proposition}[Scalar normalization]
\statusProved{}
For nonzero finite-dimensional Hilbert spaces $H,K$ and a linear map $T:H\to K$, the
map $\Phi_T(\rho)=T\rho T^*$ is completely positive and is trace-nonincreasing
exactly when $T^*T\leq1$. For every complex $c$, $\Phi_{cT}=\lvert c\rvert^2\Phi_T$.
The scalar-quotient stabilizer prescription stipulates no actual successful branch or probability. A branch can be specified using an admissible representative or an admissible class-invariant normalization rule; a phase change alone does not change the branch. Compatibility of an added normalization rule with composition is a separate obligation.
\end{proposition}
\begin{scope}
Zero outcomes are retained and never conditionally normalized. This describes the additional data for an operational lift, not a claimed canonical choice of that data.
\end{scope}
\provenance{SP-SCALAR}{Inherited: FRP-CP; reviewed structured process proofs}{Exact finite process controls}{Single blind review, sole repair and process-cluster adjudication}

\begin{proof}
For every finite-dimensional auxiliary Hilbert space $E$ and $R\geq0$,
\[
 (1_E\otimes\Phi_T)(R)=(1_E\otimes T)R(1_E\otimes T)^*\geq0,
\]
including when $T=0$, so $\Phi_T$ is completely positive.  A direct
rectangular matrix calculation gives
\[
 \operatorname{Tr}_K(T\rho T^*)=\operatorname{Tr}_H(T^*T\rho).
\]
If $T^*T\leq1$, positivity of $1-T^*T$ proves trace nonincrease.  Conversely,
applying trace nonincrease to every $\rho=|x\rangle\langle x|$ gives
$\langle x,(1-T^*T)x\rangle\geq0$, hence $T^*T\leq1$.

For every $c\in\mathbb C$,
\[
 \Phi_{cT}(\rho)=c\overline c\,T\rho T^*
 =|c|^2\Phi_T(\rho).
\]
Thus unit-modulus rescaling leaves the branch fixed.  If $T\ne0$ and
$|c|\ne1$, it changes the actual CP map and every nonzero success weight;
the zero operator remains the zero branch for every scalar.  The scalar
quotient identifies all nonzero rescalings but prescribes no actual branch,
probability, representative or normalization rule.  One may choose an
admissible representative.  Alternatively, the class-invariant rule
\[
 \mathcal N_{[T]}(\rho)=\frac{T\rho T^*}{\lVert T\rVert^2}
\]
is independent of the representative because numerator and denominator both
scale by $|c|^2$, and it is TNI because
$T^*T/\lVert T\rVert^2\leq1$.  The zero class is handled separately by the
zero branch.  Neither method is stipulated by the scalar quotient, and
compatibility of an added rule with composition is a separate obligation;
no functor or no-go theorem is asserted.  A zero-probability output is
retained without conditional normalization.
\end{proof}

\begin{proposition}[Closure and normalization of finite instruments]
\statusProved{}
For finite quantum systems with ordinary block traces, the Kraus branch prescription
is equivalent to complete positivity and trace nonincrease. Channels, branches and
instruments compose and tensor with the stated normalizations. Instruments have a
common source and target; retaining outcomes sends the output to the direct sum of
the outcome blocks, without a factor equal to the number of outcomes. Sequential
instruments retain the ordered pair of outcomes, and summing over those pairs gives a
channel. Adjoint maps under trace duality reverse the direction; a channel has a
unital adjoint.
\end{proposition}
\begin{scope}
This is an ambient process theorem. It makes no claim that arithmetic or stabilizer generators exhaust these maps.
\end{scope}
\provenance{SP-CP}{Inherited: FRP-CP; reviewed structured process proofs}{Exact finite process controls}{Single blind review, sole repair and process-cluster adjudication}

\begin{proof}
For a finite Kraus family $K_{ba,j}:H_a\to K_b$, every amplified output
block is a finite sum
\[
 \sum_{a,j}(1\otimes K_{ba,j})R_a(1\otimes K_{ba,j})^*,
\]
and is positive.  Conversely, if
$\Phi:\bigoplus_a\operatorname{End}(H_a)\to
\bigoplus_b\operatorname{End}(K_b)$ is CP, every component
$\pi_b\Phi\iota_a$ is CP.  Watrous's Theorem~2.22 supplies a finite Kraus
list for each nonzero component; a zero component receives the empty list.
Linearity reassembles these component lists into the displayed block map.

Put $A_a=\sum_{b,j}K_{ba,j}^*K_{ba,j}$.  Ordinary rectangular trace gives
\[
 \operatorname{Tr}_Y\Phi(\rho)
 =\sum_a\operatorname{Tr}_{H_a}(A_a\rho_a).
\]
Therefore trace nonincrease is equivalent to $A_a\leq1_{H_a}$ for every
input block: necessity follows by supporting a rank-one positive input on
one chosen block.  Trace preservation is similarly equivalent to
$A_a=1_{H_a}$ for every $a$.  Kraus lists are presentations of an actual CP
map; changing a list without changing that map does not change the ambient
arrow, and does not identify the finer source-certified arithmetic arrows.

For a second branch with Kraus maps $L_{cb,t}$, composition has Kraus paths
$(b,j,t)$ and operators $L_{cb,t}K_{ba,j}$.  Its effect is bounded by the
first effect because the second effects are bounded by the identity.  Tensor
has operators $K_{ba,j}\otimes M_{b'a',u}$ on Cartesian block tags.  Its
effect is $A_a\otimes D_{a'}\leq1$, and the matrix-unit spanning argument
establishes the tensor formula on entangled as well as product inputs.

For an instrument $(\Phi_o)_{o\in O}$, the retained channel has blocks
$(o,b)$ and
\[
 \widehat\Phi(\rho)_{(o,b)}=\Phi_o(\rho)_b,
 \qquad
 \operatorname{Tr}(\widehat\Phi(\rho))
 =\sum_o\operatorname{Tr}(\Phi_o(\rho))
 =\operatorname{Tr}(\rho).
\]
No outcome-count factor occurs.  Sequential instruments have branches
$\Psi_r\Phi_o$ indexed by the earlier/later pair $(o,r)$; summing over those
pairs is the composite of the two channel sums.  Independent tensor outcomes
are the listed-factor pairs $(o,s)$.

For an arbitrary complex-linear block map, choose ordinary-trace-orthonormal
matrix units $(E_p)_p$ and $(F_q)_q$ and write
\[
 \Phi(E_p)=\sum_q C_{q,p}F_q.
\]
The prescription
\[
 \Phi^{\operatorname{tr}*}(F_q)=\sum_p\overline{C_{q,p}}E_p
\]
extends linearly and satisfies
$\langle\Phi^{\operatorname{tr}*}(y),x\rangle
=\langle y,\Phi(x)\rangle$ on the matrix-unit bases and hence on every
$x,y$.  If two maps satisfy this identity, their difference is orthogonal to
every matrix unit and is zero, proving existence and uniqueness for every
linear map.  For a CP map this general adjoint specializes to
\[
 (\Phi^{\operatorname{tr}*}(y))_a
 =\sum_{b,j}K_{ba,j}^*y_bK_{ba,j}.
\]
It is CP and reverses direction.  Its value on the output identity is $A_a$,
so branches have subunital adjoints and channels have unital adjoints.
This does not make the adjoint a reverse branch: the adjoint of ordinary
discard $\operatorname{Tr}:M_2\to\mathbb C$ sends
$\lambda\mapsto\lambda I_2$ and increases the trace of $1$ from one to two.
The theorem therefore supplies ambient trace adjoints, without an automatic
CP dagger on the trace-nonincreasing branch category.
\end{proof}

\begin{proposition}[Coherent sums and tagged quantum systems]
\statusProved{}
For an odd prime $p$, the complex span of the actual stabilizer maps $\ell^2(\mathbb
F_p^n)\to\ell^2(\mathbb F_p^m)$ is the full space of linear maps. The matrix
completion therefore realizes all linear maps between its coherent Hilbert sums. For
a nonempty list of these model spaces $(H_i)_i$, the coherent algebra has dimension
$(\sum_i\dim H_i)^2$ and the tagged algebra has dimension $\sum_i(\dim H_i)^2$; the
tagged algebra is its block-diagonal subalgebra, with the corresponding
trace-preserving block-dephasing channel.
\end{proposition}
\begin{scope}
Adding coherent linear combinations enlarges the pure stabilizer fragment. A classical rig source is still to be constructed.
\end{scope}
\provenance{SP-SUM}{Inherited: F1-REAL; reviewed matrix-unit and block proof}{Exact finite sum controls}{Single blind review, sole checker repair and sum adjudication}

\begin{proof}
For $x\in\mathbb F_p^n$, tensor the computational zero preparations and
translate them to obtain an actual map
\[
 e_x:\mathbb C\longrightarrow\ell^2(\mathbb F_p^n),
 \qquad e_x(1)=\delta_x.
\]
At $n=0$ this is the identity on $\mathbb C$ indexed by the unique empty
tuple.  Adjoint closure gives $e_x^*$, so for
$y\in\mathbb F_p^m$ the actual composite
\[
 E_{y,x}=e_y e_x^*
\]
sends $\delta_z$ to $[x=z]\delta_y$.  These are the
$p^{m+n}$ rectangular computational matrix units and form a basis of all
linear maps between the two model spaces.  Individual matrix units are
actual amplitudes; their arbitrary complex sums enter at the prescribed
linear completion.  The projective scalar quotient does not supply this
addition.

This enlargement is strict.  For fixed ranks there are only finitely many
affine Lagrangian relations over the finite field, so the admitted
relation/stabilizer equivalence gives only finitely many nonzero projective
stabilizer rays.  In contrast, $\operatorname{End}(\ell^2(\mathbb F_p))$
has dimension $p^2>1$ and infinitely many rays: for independent matrix units
$E,F$, the rays of $E+zF$, $z\in\mathbb C$, are pairwise distinct.  The full
complex span therefore contains maps outside the pure actual-amplitude
fragment.

For ordered lists $X=(H_i)_{i=1}^s$ and $Y=(K_j)_{j=1}^t$, the prescribed
block action is
\[
 (T_{ji})(\xi_i)_i=\left(\sum_iT_{ji}\xi_i\right)_j.
\]
It is injective because inputs supported on one summand and projection to one
output recover every $T_{ji}$.  It is surjective because a linear map $L$
has blocks $\pi_jL\iota_i$, which reconstruct $L$ by linearity.  Matrix
composition and adjoint transpose realize ordinary composition and Hilbert
adjoint.  If either list is empty, its realization is the zero Hilbert space
and the unique empty block matrix is the unique typed zero map.  This empty
list is distinct from the one-entry list $(\mathbb C)$, and no channel on the
zero space is asserted.  Equal Hilbert spaces in repeated list positions
remain separate summands.

For a nonempty list put $H=\bigoplus_iH_i$, $d_i=\dim H_i$, and
$D=\sum_i d_i$.  Block extraction gives
\[
 \operatorname{End}(H)=\bigoplus_{i,j}\operatorname{Hom}(H_j,H_i),
\]
so its dimension is $D^2$.  The tagged algebra retains only diagonal blocks
$\bigoplus_i\operatorname{End}(H_i)$ and has dimension
$\sum_i d_i^2$.  With at least two recorded positions, any nonzero
cross-block matrix unit exhibits the possible proper inclusion, including
when the two underlying Hilbert spaces are equal.

The prescribed summand projections $P_i^X$ satisfy $P_i^X P_j^X=[i=j]P_i^X$ and $\sum_iP_i^X=1_H$.  Therefore
\[
 \Delta_X(a)=\sum_iP_i^X aP_i^X
\]
fixes each diagonal block and kills every off-diagonal block.  It is unital,
idempotent, and has exactly the tagged algebra as its range.  The projections
are a Kraus list with
\[
 \sum_i(P_i^X)^*P_i^X=1_H,
\]
so the admitted finite-block Kraus criterion makes $\Delta_X$ a channel.
In an orthonormal basis adapted to the summands,
\[
 \operatorname{Tr}_H(\Delta_X(a))
 =\sum_i\operatorname{Tr}_{H_i}(a_{ii})
 =\operatorname{Tr}_H(a),
\]
with no division by the number or dimensions of the blocks.  This channel
statement applies only to nonempty lists; no biproduct, fusion, rig,
Gaussian-closure, or source-exhaustion conclusion is added.
\end{proof}

\begin{proposition}[Fock completion and the exponential law]
\statusProved{}
On complex Hilbert spaces and contractions, the symmetric Fock prescription defines a
functor and has a natural unitary exponential law $\Gamma_s(H\oplus
K)\simeq\Gamma_s(H)\otimes\Gamma_s(K)$ with the standard normalized symmetric
tensors. In particular $\Gamma_s(0)=\mathbb C$ and $\Gamma_s(\mathbb
C)\simeq\ell^2(\mathbb N_0)$, with particle number acting diagonally on the latter.
\end{proposition}
\begin{scope}
Infinite boundedness and the completion require written arguments. No equivalence between an unspecified classical free monoid and this Hilbert construction is asserted.
\end{scope}
\provenance{SP-FOCK}{theory/symplectic-phantasm/fock.md; structured proof}{Exact finite/symbolic completion controls; scope in the DAG}{theory/verdicts/phantasm-completions-adjudication.md}

\begin{proof}[Symmetric-sector and exponential-law argument]
The permutation average is a self-adjoint idempotent, hence the orthogonal
projection onto symmetric tensors. Since $T^{\otimes r}$ commutes with every
permutation, it maps $\operatorname{Sym}^rH$ to
$\operatorname{Sym}^rK$. For a contraction,
\[
 \left\|T^{\otimes r}|_{\operatorname{Sym}^rH}\right\|
 \leq\|T\|^r\leq1.
\]
The sectorwise direct sum is therefore bounded. Its vacuum sector is the
identity, so its norm is one. Identity and composition hold sectorwise and
extend by continuity, proving functoriality. If $\|T\|>1$, choose a unit vector $h$ with $\|Th\|>1$. The symmetric
tensors $h^{\otimes r}$ have image norm $\|Th\|^r$, so no bounded
expansive-map extension is obtained.

For $r=n+m$, split $S_r$ into $(n,m)$-shuffle cosets. Internal permutations
fix $\xi_n\in\operatorname{Sym}^nH$ and
$\eta_m\in\operatorname{Sym}^mK$, while different shuffles occupy orthogonal
$H/K$ slot patterns. It follows that
\[
 \left\|P_r^{H\oplus K}
 ((\jmath_H^{\otimes n}\xi_n)\otimes
  (\jmath_K^{\otimes m}\eta_m))\right\|^2
 =\frac{n!m!}{r!}\|\xi_n\|^2\|\eta_m\|^2.
\]
The prescribed positive factor $\sqrt{r!/(n!m!)}$ makes every homogeneous
map isometric. Distinct bidegrees are orthogonal. Expanding symmetrized pure
tensors in $H\oplus K$ shows that the completed bidegree images exhaust each
sector. Thus the algebraic map extends uniquely to a unitary
\[
 \mathsf{Exp}_{H,K}:\Gamma_s(H)\otimes\Gamma_s(K)
 \longrightarrow\Gamma_s(H\oplus K).
\]
Its adjoint has the proposition's displayed direction. On homogeneous
tensors, commuting a contraction with the symmetrizer proves
\[
 \mathsf{Exp}_{H',K'}(\Gamma_s(S)\otimes\Gamma_s(T))
 =\Gamma_s(S\oplus T)\mathsf{Exp}_{H,K};
\]
density and boundedness extend this naturality equality. Vacuum and zero
factors reduce to the prescribed unit identifications.

For one mode, $e^{\otimes r}$ is a unit vector spanning the $r$-particle
sector. The polynomial vectors $x^r/\sqrt{r!}$ are orthonormal for the
existing factorial inner product, so mapping them to $e^{\otimes r}$ extends
to the one-mode completion and then to the standard basis of
$\ell^2(\mathbb N_0)$. Particle number has eigenvalue $r$ and domain
$\sum_r r^2|c_r|^2<\infty$. No Hall multiplication or CCR claim is used.
\end{proof}

\begin{proposition}[Restriction of scalars preserves the Weyl datum]
\statusProved{}
For every finite extension $E/K$ of odd-characteristic finite fields and symplectic
$E$-space $(V,\omega_E)$, the trace form on $\operatorname{Res}_{E/K}V$ is
nondegenerate and alternating. If $\chi_K$ is nontrivial, so is
$\chi_{E/K}=\chi_K\circ\operatorname{Tr}_{E/K}$, and the identity on the underlying
Weyl labels identifies the two half-form $*$-algebras exactly. These comparisons
compose along towers by transitivity of trace.
\end{proposition}
\begin{scope}
The extension degree may be divisible by the characteristic. This is restriction of scalars with a trace form, not the uncorrected inclusion of a subfield as a symplectic subsystem.
\end{scope}
\provenance{SP-TRACE}{Inherited: FRB-TRACE; structured arithmetic proof}{Exact arithmetic controls; existing bridge scope retained}{theory/verdicts/phantasm-arithmetic-adjudication.md}

\begin{proof}[Restriction-of-scalars argument]
Let $d=[E:K]$. The relative trace is $K$-linear and surjective, including
when the characteristic divides $d$. Thus
\[
 \omega_{\mathrm{Res}}(v,w)
 =\operatorname{Tr}_{E/K}(\omega_E(v,w))
\]
is alternating and $K$-bilinear. If $v\ne0$, choose $w_0$ with
$c=\omega_E(v,w_0)\ne0$ and choose $x$ with
$\operatorname{Tr}_{E/K}(x)=1$. Then
$w=(x/c)w_0$ satisfies $\omega_{\mathrm{Res}}(v,w)=1$, proving
nondegeneracy in arbitrary rank. Trace surjectivity also proves that
$\chi_{E/K}=\chi_K\circ\operatorname{Tr}_{E/K}$ is nontrivial, without
identifying the named base character with the fixed absolute-trace character.

In odd characteristic, $1/2\in K$ and
\[
 \operatorname{Tr}_{E/K}(\omega_E(v,w)/2)
 =\omega_{\mathrm{Res}}(v,w)/2.
\]
Consequently
\[
 \chi_{E/K}(\omega_E(v,w)/2)
 =\chi_K(\omega_{\mathrm{Res}}(v,w)/2),
\]
so the identity on the underlying labels preserves the half-form Weyl
product. It also preserves negation, the zero label and the coefficient
trace, and is therefore the exact unital $*$-isomorphism in the proposition.
At rank zero it is the identity on $\mathbb C$.

For a named tower $L\subset K\subset E$ and characters induced from one
named character of $L$, trace transitivity identifies the direct and
iterated forms and characters. Both comparison routes are identity on the
underlying labels and therefore compose literally. An unrelated named
character of $K$ is not silently identified with the induced one. This is
restriction of scalars through trace, rather than uncorrected subfield
inclusion or subsystem decoding.
\end{proof}

\begin{proposition}[Named arithmetic Frobenius covariance]
\statusProved{}
For a finite extension $E/K$ of odd-characteristic finite fields, with $\lvert
K\rvert=p^s$, $n\geq0$ and a named nontrivial additive character $\chi_K$, put
$\chi_{E/K}=\chi_K\circ\operatorname{Tr}_{E/K}$. On the standard register $E^n\oplus
E^n$ with the trace-form Weyl datum for $\chi_{E/K}$, coordinatewise $\#K$-power
Frobenius is $K$-linear symplectic and its specified permutation unitary satisfies
$U_{E/K,n}W^{\mathrm s}_{E,n}(a,b)U_{E/K,n}^*=W^{\mathrm
s}_{E,n}(\sigma_{E/K}(a),\sigma_{E/K}(b))$. Its induced state channel is the
invertible channel $\rho\mapsto U_{E/K,n}\rho U_{E/K,n}^*$ and its $[E:K]$-th power
is identity.
\end{proposition}
\begin{scope}
A smaller period is allowed. Abstract spaces require named semilinear data; Frobenius is not identified with partial trace.
\end{scope}
\provenance{SP-FROB}{Inherited: FRB-FROB,FRB-TRACE; channel typing: SP-CP; structured arithmetic proof}{Exact arithmetic controls; existing bridge scope retained}{theory/verdicts/phantasm-arithmetic-adjudication.md}

\begin{proof}[Relative Frobenius argument]
Put $q=|K|=p^s$, $d=[E:K]$ and $\sigma(x)=x^q$. This is a $K$-linear field
automorphism and $\sigma^d=1$. The conjugate-sum formula
\[
 \operatorname{Tr}_{E/K}(x)=\sum_{r=0}^{d-1}\sigma^r(x)
\]
shows that relative trace is invariant under $\sigma$. Hence coordinatewise
$\sigma$ preserves the restricted symplectic form and every named relative
character $\chi_K\circ\operatorname{Tr}_{E/K}$. This does not assert that an
arbitrary base character is invariant under the one-step absolute
$p$-power action.

The specified basis permutation satisfies
$U_{E/K,n}\delta_y=\delta_{\sigma y}$. On basis vectors the symmetrized Weyl
operator is
\[
 W^{\mathrm s}_{E,n}(a,b)\delta_y
 =\chi_{E/K}(-b\cdot y-a\cdot b/2)\delta_{y+a}.
\]
Conjugating by the permutation, then using character invariance and applying
$\sigma$ to the phase argument, gives
\[
 U_{E/K,n}W^{\mathrm s}_{E,n}(a,b)U_{E/K,n}^*
 =W^{\mathrm s}_{E,n}(\sigma a,\sigma b).
\]
Both phase coordinates are transformed. At $n=0$ all operators are the
identity on $\mathbb C$.

The single Kraus operator $U_{E/K,n}$ is unitary, so its conjugation is an
ordinary-trace channel. Its inverse is conjugation by the adjoint, and its
$d$-th power is identity because $U_{E/K,n}^d=1$. No minimal-period claim,
partial trace, characteristic-two half-form, or non-bijective Frobenius map
is used.
\end{proof}

\begin{proposition}[Subsystem inclusion and its dual decoder]
\statusProved{}
For an odd-characteristic finite field, a common nontrivial character and a
symplectic injection $j:U\to V$, the orthogonal decomposition $V=j(U)\oplus
j(U)^\perp$ admits the specified quantum subsystem realization. Its observable
inclusion is a unital $*$-homomorphism and its trace-dual decoder is a channel. The
decoder restricts Weyl characteristic functions to $j(U)$ and is unchanged by an
overall phase of the implementing unitary. For compatible iterated tensor
decompositions, the decoders compose as partial traces.
\end{proposition}
\begin{scope}
The complement and Weyl model compatibility must be verified. Generic field embeddings and coordinate traces are not automatically such injections.
\end{scope}
\provenance{SP-SUBSYS}{Inherited: FRP-CP; structured arithmetic proof}{Exact finite arithmetic controls}{theory/verdicts/phantasm-arithmetic-adjudication.md}

\begin{proof}[Subsystem geometry and model argument]
For a symplectic injection $j:U\to V$, the image $j(U)$ is nondegenerate,
so $j(U)\cap j(U)^\perp=0$. The dimension identity
$\dim A+\dim A^\perp=\dim V$ gives
\[
 V=j(U)\oplus j(U)^\perp.
\]
The restricted form on the complement is nondegenerate because a vector in
its radical also lies in the double perpendicular $j(U)$. Thus
$(u,w)\mapsto ju+w$ is a symplectic isomorphism, including zero retained or
zero complement factors.

The tensor of the chosen irreducible Weyl models for $U$ and the complement,
and the chosen model for $V$ transported through this symplectic
isomorphism, are irreducible unitary models of the same full matrix Weyl
algebra. Admitted model uniqueness supplies a unitary $J$ satisfying
\[
 J(W_U^{\mathrm s}(u)\otimes W_W^{\mathrm s}(w))J^*
 =W_V^{\mathrm s}(ju+w),
\]
unique up to overall phase. This reuses the admitted realization rather than
reproving finite Stone--von Neumann.

For composable $U\xrightarrow{j}V\xrightarrow{\ell}X$, the map
\[
 W_j\oplus W_\ell\longrightarrow W_{\ell j},
 \qquad(w,z)\longmapsto\ell(w)+z
\]
lands in the combined complement. Its two images are orthogonal, it is
injective, dimensions agree, and it preserves the two forms; it is therefore
symplectic. The same model comparison supplies the three registered
injection unitaries and the complement unitary. The staged and direct routes
implement the same Weyl label
$\ell(ju)+\ell(w)+z$, so unitary-intertwiner uniqueness makes them differ by
one allowed phase. This is exactly the registered pure-tensor compatibility,
including the ordinary reassociation and rank-zero unitors.
\end{proof}

\begin{proof}[Observable and decoder argument]
Conjugation gives
\[
 \iota_J(a)=J(a\otimes1)J^*,
\]
which is visibly linear, multiplicative, star-preserving and unital. For an
orthonormal basis $(e_r)$ of the complement model, put
\[
 K_r=(1\otimes\langle e_r|)J^*:H_V\to H_U.
\]
Then
\[
 \mathcal D_J(\rho)=\sum_rK_r\rho K_r^*
 =\operatorname{Tr}_{H_W}(J^*\rho J),
 \qquad \sum_rK_r^*K_r=1_{H_V}.
\]
The decoder is therefore a completely positive ordinary-trace-preserving
channel, with no division by the discarded dimension, and
\[
 \operatorname{Tr}(\mathcal D_J(\rho)a)
 =\operatorname{Tr}(\rho\iota_J(a)).
\]
Putting the complement Weyl label equal to zero gives
$\iota_J(W_U^{\mathrm s}(u))=W_V^{\mathrm s}(ju)$ and hence
\[
 \operatorname{Tr}(\mathcal D_J(\rho)W_U^{\mathrm s}(u))
 =\operatorname{Tr}(\rho W_V^{\mathrm s}(ju)).
\]
Replacing $J$ by a phase multiple changes neither the inclusion nor the
decoder.

For the registered compatible iterated data, let $A$ be the staged unitary
and $B$ the direct unitary after complement comparison and reassociation.
Compatibility says $A=\lambda B$. The phase cancels from conjugation, and
the complement unitary carries the tensor-product orthonormal basis to an
orthonormal basis of the combined complement. Expanding both partial traces
in these bases proves
\[
 \iota_{J_\ell}\iota_{J_j}=\iota_{J_{\ell j}},
 \qquad
 \mathcal D_{J_j}\mathcal D_{J_\ell}=\mathcal D_{J_{\ell j}}.
\]
This comparison is asserted only for the full registered compatibility data.
It is a trace-preserving tensor-subsystem decoder, not the trace-nonincreasing
success branch of a field support code. The full retained code decoder
includes a separate failure outcome.
\end{proof}

\begin{proposition}[Tensor assembly and product-state existence]
\statusProved{}
For the specified finite matrix factors and trace-one positive reference operators
indexed by primes, the finite tensor embeddings form an isometric unital inductive
system. Its norm completion is a unital $C^*$-algebra, and the prescribed finite
product states extend uniquely to a state on it. The GNS prescription gives its
represented von Neumann algebra with a cyclic reference vector.
\end{proposition}
\begin{scope}
No factor type, separating-vector property, inter-prime arithmetic process or Bost--Connes identification is asserted.
\end{scope}
\provenance{SP-PRIME}{theory/symplectic-phantasm/prime-tensor.md; structured proof}{Exact finite/symbolic completion controls; scope in the DAG}{theory/verdicts/phantasm-completions-adjudication.md}

\begin{proof}[Prime tensor and GNS argument]
After the fixed increasing-order permutation of tensor factors, each finite
embedding is $a\mapsto a\otimes1$. It preserves the unit, products and
adjoints. Its operator norm is unchanged: the tensor-norm upper bound gives
one inequality, and a norm-attaining vector for $a$ tensored with any unit
vector gives the reverse inequality. Slotwise identity insertion also proves
the embedding triangles, including
\[
 \iota_{Q\varnothing}(c)=c\bigotimes_{q\in Q}^{\nearrow}1_{H_q}.
\]

Represent algebraic-limit elements at finite stages. Isometry makes their
norm independent of stage, and products and adjoints may be computed after
moving representatives to a common union of prime sets. The matrix
$C^*$-identity passes to this algebraic limit. Multiplication and adjoint are
continuous, so its norm completion is a unital $C^*$-algebra with the common
finite-stage unit.

At stage $P$, put $\rho_P=\bigotimes_{p\in P}\rho_p$ and
$\varphi_P(a)=\operatorname{Tr}(\rho_Pa)$. Positivity and trace
factorization show that this is a state. Identity insertion gives
\[
 \varphi_Q(\iota_{QP}(a))
 =\varphi_P(a)\prod_{q\in Q\setminus P}\operatorname{Tr}(\rho_q)
 =\varphi_P(a).
\]
The compatible functional has norm one and extends uniquely by continuity;
approximating $y^*y$ proves positivity on the completion.

For the GNS quotient, positivity gives
$|\varphi(a^*b)|^2\leq\varphi(a^*a)\varphi(b^*b)$. The null space is a left
ideal because
\[
 \varphi((ca)^*(ca))\leq\|c\|^2\varphi(a^*a).
\]
Consequently left multiplication is well-defined and bounded by $\|c\|$ on
the quotient and extends to its Hilbert completion as a unital
$*$-representation. The class of one is cyclic because
$\pi(a)[1]=[a]$, and the represented von Neumann algebra is the double
commutant. Rank-deficient local densities may yield a cyclic but
nonseparating vector. No faithfulness of the state or representation, and no
separating or factor property, is inferred.
\end{proof}

\begin{proposition}[Arithmetic semigroup and zeta control]
\statusProved{}
In the specified represented Bost--Connes control, $\mu_m\mu_n=\mu_{mn}$,
$\mu_n^*\mu_n=1$, and $\mu_n^*e(r)\mu_n=e(nr)$, while
$\mu_ne(r)\mu_n^*=n^{-1}\sum_{ns=r}e(s)$. The prescribed time evolution leaves the
algebra invariant and satisfies $\sigma_u(\mu_n)=n^{iu}\mu_n$ and
$\sigma_u(e(r))=e(r)$. The maps $L_n$ are unital completely positive and $\nu_n$ are
corner $*$-endomorphisms. For real $b>1$, the Gibbs operator is trace class with
$\operatorname{Tr}(e^{-bH_{\log}})=\zeta(b)$ and the prescribed density has trace
one.
\end{proposition}
\begin{scope}
The Hamiltonian has logarithmic-integer eigenvalues. No zeta-zero spectrum, universal-representation faithfulness or full KMS classification is asserted.
\end{scope}
\provenance{SP-BC-CONTROL}{theory/symplectic-phantasm/bc-control.md; structured proof}{Exact finite/symbolic completion controls; scope in the DAG}{theory/verdicts/phantasm-completions-adjudication.md}

\begin{proof}[Represented arithmetic-control argument]
The basis formulas give
\[
 \mu_n^*\delta_k=
 \begin{cases}\delta_{k/n},&n\mid k,\\0,&n\nmid k,
 \end{cases}
\]
and immediately imply $\mu_m\mu_n=\mu_{mn}$ and
$\mu_n^*\mu_n=1$. They also give
$\mu_n^*e(r)\mu_n=e(nr)$. The solutions of $ns=r$ are
$s=(r+j)/n$, $0\leq j<n$. The finite root average
\[
 \frac1n\sum_{j=0}^{n-1}e^{2\pi i kj/n}
\]
is one exactly when $n\mid k$ and zero otherwise, proving the second
conjugation identity with its factor $1/n$.

Writing $q_n=\mu_n\mu_n^*$ locally, conjugation by $\mu_n$ is a
$*$-isomorphism onto
$q_nA_{\mathrm{BC}}^{\mathrm{rep}}q_n$: its inverse on the corner is compression by
$\mu_n^*$. The latter compression preserves the represented algebra, is
unital, and is completely positive by amplified quadratic-form positivity.
No multiplicativity of this compression is asserted.

The logarithmic multiplication operator is symmetric on its weighted domain.
If a vector lies in the adjoint domain, testing against every basis vector
forces the adjoint coordinates to be $(\log m)\xi_m$; square summability is
exactly the prescribed domain. Thus the operator is self-adjoint there.
Functional calculus gives
$e^{iuH_{\log}}\delta_m=m^{iu}\delta_m$, and direct computation yields
\[
 \sigma_u(\mu_n)=n^{iu}\mu_n,\qquad \sigma_u(e(r))=e(r).
\]
These formulas preserve the algebraic generator algebra. Isometric
conjugation extends them to the represented $C^*$-algebra. Continuity on
finite generator polynomials, followed by uniform norm approximation, proves
point-norm continuity.

For $b>1$, functional calculus gives
$e^{-bH_{\log}}\delta_m=m^{-b}\delta_m$. The integral comparison
\[
 \sum_{m=1}^{\infty}m^{-b}\leq1+\int_1^\infty x^{-b}\,dx
 =1+\frac1{b-1}
\]
proves trace-class convergence. The trace is the prescribed
$\zeta(b)$, so normalization gives a positive trace-one density. This is a
positive real partition sum, with no KMS, modular, factor, global, or
zeta-zero conclusion.
\end{proof}

\subsection{Decisions beyond the bootstrap}

A normalized lift of the relation semantics must retain amplitudes and
outcome probabilities. A classical additive construction must specify its
objects, morphisms, two products and distributivity before it is called a
rig quantization. Characteristic two needs its own central-extension data
before any all-prime claim. Higher gates need phase precision and a
composition law, rather than an unqualified polynomial-degree analogy.

The independent tensor assembly and the represented arithmetic control
system do not yet define their desired combination. The next definition
must supply actual arithmetic maps between the prime contributions, their
relations and a reference state. A modular comparison then requires its
analytic hypotheses and temperature convention. A spectral claim must
name its operator, domain and trace before making an arithmetic comparison.
These are decision gates for extending the graph, not unformulated theorems
assigned a proof status.
